% ==============================================================================
% CAPÍTULO 08 - TRABAJO FUTURO
% ==============================================================================

\chapter{Trabajo Futuro}
\label{cap:trabajo_futuro}
\justifying

El presente capítulo tiene como propósito describir el conjunto de actividades, componentes y desarrollos que se contemplan ejecutar como parte de la continuidad del proyecto durante la etapa de Trabajo Terminal II (TT2). 

La continuidad del proyecto se articula en torno a tres grandes ejes: la incorporación de los microservicios correspondientes a los departamentos restantes de la DDP, la consolidación e integración de los componentes inteligentes y de procesamiento documental, y el despliegue de la solución en un entorno de nube institucional.

\section{Microservicios Faltantes por Departamento}

Para Trabajo Terminal II, se desarrollarán los microservicios pendientes de las áreas operativas de la DDP. Cada departamento tendrá un microservicio propio conectado a través del API Gateway; esto garantiza la separación de responsabilidades, permite hacer modificaciones independientes en los procesos de cada área y facilita el escalamiento autónomo según su carga de trabajo real. 

En la Tabla \ref{tab:microservicios_tt2} se detallan estos componentes, su descripción y su estado actual.

\begin{table}[H]
\centering
\caption{Microservicios por departamento contemplados para TT2}
\label{tab:microservicios_tt2}
\renewcommand{\arraystretch}{1.4}
\begin{tabularx}{\textwidth}{|p{5cm}|X|c|}
\hline
\textbf{Microservicio} & \textbf{Descripción} & \textbf{Estado} \\
\hline

Microservicio del Quejoso &
Gestiona el registro, consulta y seguimiento de quejas por parte del ciudadano. &
Implementado \\
\hline

Microservicio de Recepción &
Administra la recepción inicial de expedientes, validación de información y asignación al área correspondiente. &
Pendiente \\
\hline

Microservicio de Primer Contacto &
Permite el análisis preliminar del caso, determinación de competencia y orientación al usuario. &
Pendiente \\
\hline

Microservicio de Subdefensoría &
Gestiona la investigación interna, comunicaciones oficiales y seguimiento de oficios. &
Pendiente \\
\hline

Microservicio de la Defensoría &
Administra acuerdos de conciliación, firmas institucionales y resoluciones finales. &
Pendiente \\
\hline

Microservicio de Administración Técnica &
Centraliza la configuración, parámetros del sistema, catálogos institucionales y administración de usuarios. &
Pendiente \\
\hline

\end{tabularx}
\end{table}

La construcción de cada uno de estos microservicios contempla las etapas de análisis de requerimientos específicos por área, diseño detallado del modelo de datos, implementación de la lógica de negocio, definición de los contratos de comunicación con el API Gateway, pruebas unitarias e integración con el expediente digital centralizado.

\section{Integración de Componentes Especializados}

Durante Trabajo Terminal II se completará el desarrollo y la integración de los componentes especializados iniciados en la primera etapa. En la Tabla \ref{tab:componentes_especializados} se detallan las actividades planificadas para cada uno de ellos y su propósito en el sistema.

\begin{table}[ht]
\centering
\caption{Componentes especializados y actividades planificadas para TT2.}
\label{tab:componentes_especializados}
\begin{tabular}{|p{3.5cm}|p{5.5cm}|p{5.5cm}|}
\hline
\textbf{Componente} & \textbf{Actividades para TT2} & \textbf{Propósito / Impacto} \\ \hline
Clasificador Inteligente & Seguir entrenando el modelo con un corpus ampliado, evaluarlo con casos reales anonimizados y conectarlo al API Gateway. & Optimizar la precisión al sugerir la competencia de las quejas y realizar su integración con el departamento de Primer Contacto. \\ \hline
Compresor y Procesador Multimedia & Consolidar el procesamiento de archivos y exponerlo como un servicio independiente para los módulos que manejan documentos. & Reducir el tamaño de las evidencias digitales para ahorrar almacenamiento y ancho de banda, sin perder su valor legal. \\ \hline
Expediente Digital Centralizado & Completar la estructura del expediente para unificar datos, evidencias y resoluciones en un solo punto accesible por servicios autorizados. & Asegurar el seguimiento y la trazabilidad de la queja durante todo su ciclo de vida institucional. \\ \hline
\end{tabular}
\end{table}
\section{Despliegue en la Nube y Operación en Producción}

En Trabajo Terminal II se realizará la migración de la plataforma del entorno local a un entorno de nube para iniciar su operación institucional. Este despliegue incluirá la integración de todos los componentes del sistema, junto con la configuración de respaldos y políticas de seguridad para proteger la información sensible. 

La selección del proveedor de nube se definirá durante la segunda etapa del proyecto, tras concluir un análisis comparativo enfocado en los requisitos de procesamiento del clasificador inteligente y el almacenamiento de evidencias multimedia. Las actividades principales planificadas para esta fase se detallan en la Tabla \ref{tab:despliegue_nube}.

\begin{table}[H]
\centering
\caption{Actividades contempladas para el despliegue en la nube}
\label{tab:despliegue_nube}
\renewcommand{\arraystretch}{1.4}
\begin{tabularx}{\textwidth}{|p{5.5cm}|X|}
\hline
\textbf{Actividad} & \textbf{Descripción} \\
\hline

Selección del proveedor de nube &
Evaluación de alternativas viables conforme a las restricciones institucionales, considerando costos, soberanía de datos y disponibilidad. \\
\hline

Configuración del servidor &
Provisionamiento de instancias, redes, almacenamiento persistente y reglas de seguridad para la operación de los microservicios. \\
\hline

Contenedorización de servicios &
Empaquetado de cada microservicio, del clasificador y del compresor de archivos en contenedores para facilitar su despliegue y portabilidad. \\
\hline

Bases de datos en la nube &
Migración del esquema de datos y configuración de respaldos, replicación y políticas de retención. \\
\hline

Pruebas de carga y resiliencia &
Validación del comportamiento del sistema bajo escenarios de uso concurrente y condiciones adversas. \\
\hline

\end{tabularx}
\end{table}

\section{Consideraciones Finales}

El trabajo programado para TT2 representa el desarrollo, integración y puesta en operación de la plataforma empezada durante TT1. La estrategia adoptada sobre un microservicio por departamento nos permite que el proyecto evolucione de manera ordenada hacia su versión final, incorporando los microservicios pendientes, completando los componentes inteligentes y trasladando la solución a un entorno productivo en la nube.

De esta forma, la siguiente etapa del proyecto tiene planteado garantizar que la plataforma pueda ser realmente de utilidad para la institución, fortaleciendo los procesos internos de la DDP y consolidando una base tecnológica sostenible para su operación a futuro.
